\documentclass{article}
\usepackage[UTF8]{ctex}
\usepackage{amsmath,amssymb,amsfonts}
\usepackage{enumitem}
\usepackage{geometry}
\geometry{a4paper, margin=1in}
% === 错题本排版优化（不参与题目标识） ===
\setlength{\parindent}{0pt}
\setlength{\parskip}{0.35em plus 0.10em minus 0.08em}
\setlength{\emergencystretch}{3em}
\setlist{itemsep=0.25em, topsep=0.40em, parsep=0pt}
\allowdisplaybreaks[2]
% === 错题本排版优化结束 ===
\begin{document}
\section*{北京工业大学2026年数学分析真题}
\section*{10．北京工业大学2026年数学分析试卷}

\subsection{题目1}
1．应用柯西准则证明 \(\left\{1+\frac{1}{1!}+\frac{1}{2!}+\cdots+\frac{1}{n!}\right\}\) 收敛并求极限。

\subsection{题目2}
2．设 \(f(x)\) 在 \([-1,1]\) 上具有连续三阶导数，且 \(f(-1)=0, f(1)=1, f^{\prime}(0)=0\) ，证明：在 \((-1,1)\) 内至少存在一点 \(\xi\) ，使 \(f^{\prime \prime \prime}(\xi)=3\) ．

\subsection{题目3}
3．叙述致密性定理并证明：设 \(D_{1}, D_{2}\) 是 \(\mathbb{R}^{2}\) 的两个非空闭集且 \(D_{1}\) 有界，则存在 \(x_{0} \in D_{1}, y_{0} \in D_{2}\) ，使

\[
\left\|x_{0}-y_{0}\right\|=\inf _{\substack{x \in D_{1} \\ y \in D_{2}}}\|x-y\|
\]

其中 \(\|\cdot\|\) 表示 \(\mathbb{R}^{2}\) 中的欧氏距离。

\subsection{题目4}
4．求幂级数 \(\sum_{n=1}^{\infty}\left(1+\frac{1}{2}+\frac{1}{3}+\cdots+\frac{1}{n}\right) x^{n}\) 的收敛域与和函数。

\subsection{题目5}
5．设 \(f(x, y)=\left\{\begin{array}{ll}\frac{x y\left(x^{2}-y^{2}\right)}{x^{2}+y^{2}}, & x^{2}+y^{2} \neq 0, \\ 0, & x^{2}+y^{2}=0 .\end{array}\right.\) 求 \(f_{x y}^{\prime \prime}(0,0)\) 和 \(f_{y x}^{\prime \prime}(0,0)\) ．

\subsection{题目6}
6．证明 \(\sum_{n=0}^{\infty} x^{n}(1-x)^{2}\) 在 \([0,1]\) 上一致收敛，但 \(\sum_{n=0}^{\infty} x^{n}(1-x)\) 在 \([0,1]\) 上非一致收敛。

\subsection{题目7}
7．设抛物线 \(y=a x^{2}+b x\) 满足当 \(x \in[0,1]\) 时，\(y \geq 0\) ，已知改抛物线与 \(x\) 轴及直线 \(x=1\) 所围图形的面积为 \(\frac{1}{3}\) ，试确定 \(a, b\) 使该抛物线与直线 \(x=1\) 及 \(x\) 轴所围图形绕 \(x\) 轴旋转一周而成的旋转体的体积最小。

\subsection{题目8}
8 ．利用三重积分计算椭球体 \(\Omega:(x+y)^{2}+(y+z)^{2}+(z+x)^{2} \leq 1\) 的体积．

\subsection{题目9}
9．设一致连续的非负函数 \(f(x)\) 满足 \(\int_{1}^{+\infty} f(x) \mathrm{d} x<+\infty\) ，证明： \(\lim _{x \rightarrow+\infty} f(x)=0\) ．若将一致连续改为连续，结论不成立，举出反例．

\subsection{题目10}
10．计算 \(I=\iint_{\sum} \frac{2 \mathrm{~d} y \mathrm{~d} z}{x \cos ^{2} x}+\frac{\mathrm{d} z \mathrm{~d} x}{\cos ^{2} y}-\frac{\mathrm{d} x \mathrm{~d} y}{z \cos ^{2} z}, \sum: x^{2}+y^{2}+z^{2}=1\) ，取外侧．
\end{document}
